\answer
\begin{enumerate}

%%--- 章末問題 7.1 ---
\item[\ref{ex7.1}]
出力電力は$P_{\mathrm{out}} = 5\times0.5 = 2.5$\,Wで一定。式\ref{eq7.2}・式\ref{eq7.3}より
\begin{center}
\begin{tabular}{c|c|c}
\hline
$V_{\mathrm{in}}$ & $P_{\mathrm{loss}}=(V_{\mathrm{in}}-5)\times0.5$ & $\eta=5/V_{\mathrm{in}}$ \\
\hline
$6$\,V & $0.5$\,W & $83\,\%$ \\
$9$\,V & $2.0$\,W & $56\,\%$ \\
$12$\,V & $3.5$\,W & $42\,\%$ \\
\hline
\end{tabular}
\end{center}
入力電圧が高いほど損失が増え効率が下がる。
効率よく使うには\textbf{入力電圧を出力電圧にできるだけ近づける}（電圧差を小さくする）ことである。

%%--- 章末問題 7.2 ---
\item[\ref{ex7.2}]
(a) 式\ref{eq7.5}より
\begin{equation*}
R_s = \frac{9 - 4.7}{4\times10^{-3}} = \frac{4.3}{4\times10^{-3}} \approx 1.1\,\mathrm{k}\Omega
\end{equation*}

(b) $V_Z$が$4.7$\,Vのまま一定なら，入力が$12$\,Vに上がったときのツェナー電流は
\begin{equation*}
I_Z = \frac{12 - 4.7}{1075} \approx 6.8\,\mathrm{mA}
\end{equation*}
入力が$3$\,V上がった分は，ほぼすべて$R_s$の電圧降下（電流の増加）が引き受け，
基準電圧$V_Z$はほとんど動かない。
入力電圧がふらついても基準が一定に保たれるので，これを物差しにした出力も安定する ―
これが定電圧源として都合のよい理由である。

%%--- 章末問題 7.3 ---
\item[\ref{ex7.3}]
(a) 式\ref{eq7.9}より
\begin{equation*}
\frac{R_1}{R_2} = \frac{V_{\mathrm{out}}}{V_{\mathrm{ref}}} - 1
= \frac{12}{2.5} - 1 = 4.8 - 1 = 3.8
\end{equation*}

(b) $R_1 = 3.8 \times R_2 = 3.8 \times 4.7\,\mathrm{k}\Omega \approx 18\,\mathrm{k}\Omega$。

(c) 分圧より$V_{\mathrm{fb}} = \dfrac{R_2}{R_1+R_2}V_{\mathrm{out}}$（式\ref{eq7.8}）。
負帰還が効くと仮想短絡（式\ref{eq7.7}）で$V_{\mathrm{fb}} = V_{\mathrm{ref}}$になるから，
\begin{equation*}
V_{\mathrm{ref}} = \frac{R_2}{R_1+R_2}V_{\mathrm{out}}
\end{equation*}
これを$V_{\mathrm{out}}$について解くと
$V_{\mathrm{out}} = \dfrac{R_1+R_2}{R_2}V_{\mathrm{ref}} = \left(1+\dfrac{R_1}{R_2}\right)V_{\mathrm{ref}}$
（式\ref{eq7.9}）を得る。

%%--- 章末問題 7.4 ---
\item[\ref{ex7.4}]
(a) $P_{\mathrm{out}} = 3.3\times1.5 = 4.95$\,W。式\ref{eq7.2}より
\begin{equation*}
P_{\mathrm{loss}} = (12 - 3.3)\times1.5 = 8.7\times1.5 \approx 13\,\mathrm{W}
\end{equation*}
効率は$\eta = 3.3/12 \approx 0.28 = 28\,\%$。

(b) 降圧チョッパでは入力$P_{\mathrm{in}} = 4.95/0.88 \approx 5.6$\,W，
損失は$5.6 - 4.95 \approx 0.68$\,Wで済む。
リニアレギュレータは出力の$2.6$倍もの$13$\,Wを\textbf{パストランジスタ1個}で発熱させ，
大きな放熱器が必要になる（\ref{sec7.4}節の注意）。
電圧差が大きく電流も大きいこの用途では，損失が20分の1で済む\textbf{降圧チョッパがふさわしい}。

%%--- 章末問題 7.5 ---
\item[\ref{ex7.5}]
(a) ツェナー基準だけの回路では，出力は$V_{\mathrm{out}} = V_{\mathrm{ref}} - V_{\mathrm{BE}}$（式\ref{eq7.10}）。
負荷を軽くしていくと出力電圧がわずかに動く。
これは負荷電流の変化で$V_{\mathrm{BE}}$が変わり，それがそのまま出力に残るためである
（ツェナー電圧$V_Z$自体はほとんど動かない）。

(b) オペアンプで負帰還をかけた回路では，$V_{\mathrm{BE}}$の変動もループの内側に閉じ込められ，
同じ負荷変動に対する出力の変動は(a)より大幅に小さくなる。
定常出力は式\ref{eq7.9}の$V_{\mathrm{out}}=(1+R_1/R_2)V_{\mathrm{ref}}$とよく一致する。

(c) 入力を下げていくと，ある電圧（$V_{\mathrm{out}}+V_{\mathrm{dropout}}$，式\ref{eq7.11}）を割り込んだところで
パストランジスタが開ききって負帰還が効かなくなり，出力が入力に引きずられて低下し始める。
この崩れ始める入力電圧からドロップアウト電圧を読み取れる。

\end{enumerate}
